\chapter{Experimental Setup}
\label{chap:setup}

\section{Benchmark and Protocol}
We use all 101 WebGen-Bench tasks. They come from 20 application domains of unequal size, from eight tasks each (company brochures, learning platforms, discussion forums, e-commerce) down to two (personal blogs, streaming, email clients). The benchmark's manifest order groups tasks into domain blocks; \Cref{sec:order-design} adds shuffled orders. We follow the official protocol: the generator produces the project once, there is no repair loop, and a project is passed to evaluation as soon as its dependencies install and the development server starts. Test cases without an evaluation record count as failed, as in the official scoring script. An earlier phase of this project used a stricter build gate (\texttt{vue-tsc}) and a more generous repair loop; those results are not comparable with the paper and are not reported as main results.

\Cref{fig:pipeline} shows one task in the pipeline. The memory module is the only component that differs between arms; the generator, the scaffold, the tester and the scoring are shared.

\begin{figure}[ht]
\centering
\begin{tikzpicture}[node distance=6mm and 9mm, every node/.style={font=\small},
  box/.style={draw, rounded corners, align=center, minimum height=10mm, inner sep=4pt},
  mem/.style={box, draw=blue!70!black, fill=blue!6}, >=stealth]
\node[box] (req) {Requirement\\(task $t$)};
\node[box, right=of req] (gen) {Generator\\WebGen-LM-7B\\in Bolt.diy};
\node[box, right=of gen] (site) {Website\\(started once)};
\node[box, right=of site] (eval) {Tester and grader\\Qwen2.5-VL-32B};
\node[mem, below=12mm of gen, xshift=18mm] (bank) {Memory bank: checklist, code pool, lessons\\Qwen2.5-VL-72B + bge-m3};
\draw[->] (req) -- (gen); \draw[->] (gen) -- (site); \draw[->] (site) -- (eval);
\draw[->, blue!70!black] (eval.south) |- node[pos=0.25, right, align=left, font=\scriptsize]{reflect\\after task $t$} (bank.east);
\draw[->, blue!70!black] (bank.west) -| node[pos=0.75, left, align=right, font=\scriptsize]{memory for\\task $t{+}1$} (req.south);
\end{tikzpicture}
\caption{One task in the pipeline. Black: the official WebGen-Bench setup (No Memory). Blue: the memory module, which reflects after each task and adds text to the prompt of later tasks only.}
\label{fig:pipeline}
\end{figure}

\section{Models and Roles}
All models are self-hosted on the university CS ML cluster (two NVIDIA H200 GPUs per job, served with vLLM). No commercial \gls{llm} service is called. \Cref{tab:roles} lists the roles.

\begin{table}[ht]
\centering
\caption{Model roles. The appearance grader is the only deviation from the paper, which uses GPT-4o.}
\label{tab:roles}
\begin{tabularx}{\textwidth}{lX}
\toprule
Role & Model \\
\midrule
Generator & WebGen-LM-7B (BF16), Bolt.diy scaffold, single generation \\
Functional tester & Qwen2.5-VL-32B-Instruct (WebVoyager-style \gls{gui} agent) \\
Appearance grader & Qwen2.5-VL-32B-Instruct (paper: GPT-4o) \\
Memory / reflection & Qwen2.5-VL-72B-Instruct-AWQ \\
Embeddings & bge-m3 \\
\bottomrule
\end{tabularx}
\end{table}

The memory model is deliberately different from the evaluator, so that reflection cannot copy the evaluator's grading behaviour. vLLM applies the model repository's generation configuration by default; WebGen-LM-7B ships with a repetition penalty that Bolt.diy never sends, so we disable it for the generator to match the paper's setting.

\section{Metrics}
For each task with $n$ test cases, of which $p$ pass and $q$ are partial,
\begin{equation}
\text{accuracy} = \frac{p + 0.5\,q}{n}.
\end{equation}
We report accuracy over all test cases, the mean appearance score (1--5), and the number of \emph{zero-score} tasks. Because test cases within a task are not independent, comparisons are paired by task rather than by test case.

\section{Experimental Design}
\label{sec:design}
\Cref{tab:design} lists every arm, the task orders in which it was run and the question it answers. All arms use WebGen-LM-7B with Bolt.diy under the official protocol, cover all 101 tasks and start from an empty memory; each arm is run once per order. The task orders are defined in \Cref{sec:order-design}.

\begin{table}[ht]
\centering
\caption{Experimental design. Memory content is what is added to the generator's prompt; version numbers refer to the design history in \Cref{sec:history}.}
\label{tab:design}
\small
\begin{tabularx}{\textwidth}{>{\raggedright}p{2.9cm}>{\raggedright}X>{\raggedright}p{2.6cm}>{\raggedright\arraybackslash}X}
\toprule
Arm & Memory content & Task orders & Question \\
\midrule
No Memory & none (requirement only) & manifest, seeds 11, 29, 47, interleaved & Does the baseline reproduce the paper, and how much does it vary between runs? \\
Text memory (v4) & checklist with textual lessons & manifest & Does textual memory help a small generator? \\
Checklist only & checklist; memory neither read nor written & manifest & How much does decomposing the requirement add? \\
Code memory (v7) & checklist, lessons and code examples & manifest, seeds 11, 29 & Does code memory help, and does it survive a change of order? \\
Fixed code prefix & checklist and the final code memories from task 1; no memory & manifest & Is the gain learning, or just code in the prompt? \\
Final (v9.1) & as v7, code examples are verified fragments only & manifest, seeds 11, 29, 47, interleaved & Does the final method hold across task orders? \\
\bottomrule
\end{tabularx}
\end{table}

Memory never uses information from the current task itself (leave-one-task-out), and evaluation results are never fed back into the same task.

\section{Reproducing the Baseline}
No Memory on all 101 tasks reaches 28.5\% accuracy (paper: 28.4\%) with a mean appearance score of 2.35 (paper: 2.5 with GPT-4o) and 30 zero-score tasks. The close match in accuracy gives us confidence that the self-hosted pipeline is faithful to the official one.

\section{Reproducibility and Constraints}
Every run starts from an empty memory and stores its configuration file, task order, prompts, retrieved and injected memory, raw evaluator transcripts and scores; the key files of every run are archived in the project repository. The generator is decoded deterministically, so differences between arms come from the prompt rather than from sampling, although the \gls{gui} tester is not deterministic.

The main constraint is compute. A task takes about twelve minutes with memory and six without, so one 101-task run takes about a day on our two GPUs. As in the original benchmark, each configuration is therefore run once per task order. We compensate by pairing comparisons by task, by repeating the No Memory baseline in three orders to measure its run-to-run variation, and by repeating the memory method across orders rather than within one.
